% !PRINTED 180
\weberpage{180}
\weq{3}{\Sigma \xi_\nu u_\nu = \Sigma \xi'_\nu u'_\nu,}
also, wie in \S.\,55, (3)
\weq{4}{u'_k = \overset{i}{\Sigma} \alpha_i^{(k)} u_i.}
Nun ist
\weq{5}{u_k = \overset{i}{\Sigma} a_{k,i} x_i, \qquad u'_k = \overset{i}{\Sigma} a'_{k,i} x'_i,}
und demnach giebt uns (4) die Transformation von $m$ linearen
Functionen. Die Co\"efficienten des gegebenen Systems sind $a'_{k,h}$
und des transformirten, die man durch Einsetzen des ersten
Systems (5) in (4) erh\"alt,
\weq{6}{b_{k,h} = \overset{i}{\Sigma} \alpha_i^{(k)} a_{i,h}.}
Die Determinante des ersteren Systems
\[
H' = \Sigma \pm a'_{1,1} a'_{2,2} \ldots a'_{m,m}
\]
ist nach \S.\,55, (4)
\[
= r \Sigma \pm b_{1,1} b_{2,2} \ldots b_{m,m},
\]
und nach (6) ist die Determinante der $b_{k,h}$
\[
= r \Sigma \pm a_{1,1} a_{2,2} \ldots a_{m,m}.
\]
Wenn wir also
\weq{7}{H = \Sigma \pm a_{1,1} a_{2,2} \ldots a_{m,m}}
setzen, so haben wir den wichtigen Satz
\weq{8}{H' = r^2 H.}
Die Determinante $H$ wird die \gesperrt{Determinante der Form}
$\varphi(x)$ genannt.

Sie \"andert sich also bei der linearen Transformation der
Function $\varphi$ nur um den Factor $r^2$, und sie kann bei der
transformirten Form nicht verschwinden, wenn sie es bei der
urspr\"unglichen nicht thut.

\gesperrt{Das Verschwinden der Determinante $H$ hat die Bedeutung,
dass die Function $\varphi(x)$ sich auffassen l\"asst als
eine Function von weniger als $m$ Ver\"anderlichen, die
lineare Functionen der $x$ sind.}

Denn wenn erstens durch eine lineare Transformation $\varphi(x)$
in die Function $\Phi(x')$ \"ubergeht, die von einer der Ver\"anderlichen,
etwa von $x'_1$, unabh\"angig ist, so verschwindet $H'$, da die Elemente
in einer der Zeilen (und Colonnen) alle verschwinden, und also
ist wegen (8) auch $H = 0$.

Zweitens aber ist $H$ die Determinante des Systems linearer
Gleichungen
%CONT>>
